% 1.5
% Ficheiro incluído por main.tex

\section[Partilha SRII: SMB e permissões NTFS]{Partilha SRII: SMB e permissões NTFS}

\esquema{1-5-partilha-srii}{Partilha SRII e permissões por grupo}

\textbf{Objetivo:} criar a pasta partilhada \textbf{SRII} onde o Grupo-Alunos só lê e o Grupo-Professores lê e escreve.

\passo{Passo 1 - Criar a pasta e partilhá-la}

\begin{enumerate}
\item Criar uma pasta no \textbf{Ambiente de Trabalho (Desktop)} com o nome \textbf{\texttt{SRII}}.
\item Botão direito → \textbf{Properties} → separador \textbf{Sharing} → \textbf{Advanced Sharing...}
\item Marcar \textbf{Share this folder} → \textbf{Permissions}.
\end{enumerate}

\passo{Passo 2 - Permissões de partilha (grupos)}

\begin{enumerate}
\item \textbf{Remover} o grupo \texttt{Everyone} (por segurança).
\item \textbf{Add...} → \texttt{Grupo-Professores} → OK → coluna \emph{Allow}: \textbf{Full Control}.
\item \textbf{Add...} → \texttt{Grupo-Alunos} → OK → coluna \emph{Allow}: \textbf{apenas Read}.
\item \textbf{OK} duas vezes para voltar às propriedades.
\end{enumerate}

\begin{simplesbox}
\textbf{[VERIFICAR]} Para o Grupo-Professores, \emph{Change} também seria suficiente (com NTFS \emph{Modify} o professor já lê e escreve). Decidir qual das duas usar.
\end{simplesbox}

\passo{Passo 3 - Permissões de segurança (NTFS)}

\begin{enumerate}
\item Separador \textbf{Security} → \textbf{Edit...} → \textbf{Add...} → adicionar \texttt{Grupo-Alunos} e \texttt{Grupo-Professores}.
\item Selecionar \texttt{Grupo-Alunos} e verificar que estão marcadas \textbf{Read \& execute}, \textbf{List folder contents} e \textbf{Read}.
\item Selecionar \texttt{Grupo-Professores} e marcar \textbf{Modify} na coluna \emph{Allow}.
\item \textbf{OK} para aplicar e fechar as propriedades.
\end{enumerate}

\passo{Passo 4 - Ficheiro de teste}

Criar dentro da SRII um ficheiro de texto \textbf{\texttt{ficheiroteste.txt}} (com qualquer conteúdo), para os testes de leitura no Ubuntu.

\begin{simplesbox}
\textbf{[VERIFICAR]} Nome do ficheiro de teste: \texttt{ficheiroteste.txt} ou \texttt{notas.txt} - usar o mesmo em todo o workshop.
\end{simplesbox}

\emph{Verificação esperada:} a pasta mostra o ícone de partilha e contém \texttt{ficheiroteste.txt}.

\problemas{\hyperref[ec-a]{Erros Comuns - Parte A}}
